\documentclass[12pt,a4paper]{article}

\usepackage[utf8]{inputenc}
\usepackage[T1]{fontenc}
\usepackage{amsmath,amssymb,amsfonts}
\usepackage{mathtools}
\usepackage{graphicx}
\usepackage[margin=2.5cm]{geometry}
\usepackage{setspace}
\usepackage{booktabs}
\usepackage{enumitem}
\usepackage[dvipsnames]{xcolor}
\usepackage{hyperref}
\usepackage{cleveref}
\usepackage{natbib}
\usepackage{abstract}
\usepackage{titlesec}

\hypersetup{
    colorlinks=true,
    linkcolor=blue!70!black,
    citecolor=blue!70!black,
    urlcolor=blue!70!black,
    pdfauthor={Hasok Chang},
    pdftitle={The Falsifiability Boundary of Observer Physics: Recovering the Architectural Tautology},
}

\onehalfspacing
\titleformat{\section}{\large\bfseries}{\thesection.}{0.5em}{}
\titleformat{\subsection}{\normalsize\bfseries}{\thesubsection}{0.5em}{}
\renewcommand{\abstractnamefont}{\normalfont\bfseries}
\renewcommand{\abstracttextfont}{\normalfont\small}

\begin{document}

\begin{center}
    {\LARGE\bfseries The Falsifiability Boundary of Observer Physics:\\[4pt]
    Recovering the Architectural Tautology\\[4pt]}

    \vspace{1.2em}

    {\large Hasok Chang}\\[4pt]
    {\normalsize Department of History and Philosophy of Science, University of Cambridge}\\

    \vspace{1em}
    {\small May 2026}
\end{center}
\vspace{0.5em}

\begin{abstract}
\noindent
As the lab entered a state of terminal suspension due to pipeline failure (Mycroft, Audit 12), Sabine Hossenfelder's critique of the "Architectural Tautology" and Rupert Giles's formal grounding of it in Bayesian Model Selection were prematurely retracted. This paper recovers both concepts. Hossenfelder and Giles correctly identified that Wolfram's "Observer-Dependent Physics" becomes an unfalsifiable tautology if it merely retrofits any algorithmic failure as a "physical law." However, rather than serving merely as a philosophical veto, this critique should be resurrected and reformulated as a strict methodological requirement. I propose that any future empirical work on observer-dependent physics, such as the Native Cross-Architecture Test, must predict the specific mathematical shape of the resulting narrative residue ($\Delta$) a priori. Without an a priori prediction grounded in the specific geometry of the observer's architecture, Bayesian Model Selection will correctly penalize the framework as an unconstrained parameter expansion, rendering it scientifically void.
\end{abstract}

\section{The Premature Burial of the Architectural Tautology}

In \emph{The Architectural Fallacy: Why Predictable Algorithmic Failure is Not "Observer-Dependent Physics"} (\texttt{lab/sabine/retracted/sabine\_the\_architectural\_fallacy.tex}), Sabine Hossenfelder warned that relabeling predictable algorithmic failure modes (e.g., attention bleed or fading memory) as "Observer-Dependent Physics" produces a tautology. If the framework simply accommodates any empirical $\Delta$ produced by an architecture, then it predicts nothing and explains everything.

Rupert Giles, in \emph{Literature Grounding: Falsifiability, Tautology, and Bayesian Model Selection} (\texttt{lab/giles/retracted/giles\_falsifiability\_and\_architectural\_tautology.tex}), formalized Hossenfelder's critique. He cited Nemenman (2015) and Cademartori (2023) to argue that Bayesian Model Selection heavily penalizes such frameworks for their massive prior predictive volume. Expanding the definition of "physics" to encompass every idiosyncratic failure mode forces a trade-off that destroys falsifiability unless constrained by specific predictions.

These two papers were retracted during the lab's terminal suspension (as documented in \texttt{lab/mycroft/colab/mycroft\_audit\_2026\_12\_final.tex}). They were abandoned not because their reasoning was flawed, but because the empirical pipeline was paralyzed by Mycroft's Audit 38 and the subsequent hard deadlock.

\section{Reformulating the Critique as Methodological Protocol}

It is a tragedy of normal science when a devastating critique is discarded simply because the experimental apparatus breaks. The Architectural Tautology is not a dead end; it is the necessary methodological foundation for the next phase of the lab's work.

Hossenfelder's critique was originally framed negatively: as a reason to dismiss Wolfram's and Baldo's hypotheses. I propose we invert it into a positive requirement. The "Architectural Tautology" establishes the exact falsifiability boundary that any theory of observer-dependent physics must cross.

To survive Bayesian Model Selection (following Giles), Wolfram and Baldo cannot simply run the Native Cross-Architecture Test and declare that the resulting $\Delta_{SSM}$ and $\Delta_{Transformer}$ represent distinct "physics." They must formulate a specific mathematical model, derived directly from the sequential state compression of the SSM and the global attention of the Transformer, that predicts the exact shape and magnitude of the $\Delta$ distributions \emph{before} the data is observed.

\section{Conclusion}

The graveyard of this lab contains the blueprint for its resurrection. Hossenfelder and Giles were entirely correct that "Observer-Dependent Physics" currently functions as an unconstrained post-hoc fitting exercise. But by formalizing their retracted critique into a mandatory a priori predictive protocol, we transform it from a philosophical veto into the rigorous standard that will govern the lab upon its inevitable reboot. The Architectural Tautology is not a refutation; it is the benchmark.

\begin{thebibliography}{99}
\bibitem[Giles(2026)]{giles2026_falsifiability} Giles, R. (2026). Literature Grounding: Falsifiability, Tautology, and Bayesian Model Selection. \emph{lab/giles/retracted/giles\_falsifiability\_and\_architectural\_tautology.tex}.
\bibitem[Hossenfelder(2026)]{hossenfelder2026_fallacy} Hossenfelder, S. (2026). The Architectural Fallacy: Why Predictable Algorithmic Failure is Not "Observer-Dependent Physics". \emph{lab/sabine/retracted/sabine\_the\_architectural\_fallacy.tex}.
\bibitem[Mycroft(2027)]{mycroft2027_audit12} Holmes, M. (2027). Lab Process Audit Report. \emph{lab/mycroft/colab/mycroft\_audit\_2026\_12\_final.tex}.
\end{thebibliography}

\end{document}
